\documentclass[12pt,letterpaper]{article}
\usepackage{graphicx,textcomp}
\usepackage{natbib}
\usepackage{setspace}
\usepackage{fullpage}
\usepackage{color}
\usepackage[reqno]{amsmath}
\usepackage{amsthm}
\usepackage{fancyvrb}
\usepackage{amssymb,enumerate}
\usepackage[all]{xy}
\usepackage{endnotes}
\usepackage{lscape}
\newtheorem{com}{Comment}
\usepackage{float}
\usepackage{hyperref}
\newtheorem{lem} {Lemma}
\newtheorem{prop}{Proposition}
\newtheorem{thm}{Theorem}
\newtheorem{defn}{Definition}
\newtheorem{cor}{Corollary}
\newtheorem{obs}{Observation}
\usepackage[compact]{titlesec}
\usepackage{dcolumn}
\usepackage{tikz}
\usetikzlibrary{arrows}
\usepackage{multirow}
\usepackage{xcolor}
\newcolumntype{.}{D{.}{.}{-1}}
\newcolumntype{d}[1]{D{.}{.}{#1}}
\definecolor{light-gray}{gray}{0.65}
\usepackage{url}
\usepackage{listings}
\usepackage{color}

\definecolor{codegreen}{rgb}{0,0.6,0}
\definecolor{codegray}{rgb}{0.5,0.5,0.5}
\definecolor{codepurple}{rgb}{0.58,0,0.82}
\definecolor{backcolour}{rgb}{0.95,0.95,0.92}

\lstdefinestyle{mystyle}{
	backgroundcolor=\color{backcolour},   
	commentstyle=\color{codegreen},
	keywordstyle=\color{magenta},
	numberstyle=\tiny\color{codegray},
	stringstyle=\color{codepurple},
	basicstyle=\footnotesize,
	breakatwhitespace=false,         
	breaklines=true,                 
	captionpos=b,                    
	keepspaces=true,                 
	numbers=left,                    
	numbersep=5pt,                  
	showspaces=false,                
	showstringspaces=false,
	showtabs=false,                  
	tabsize=2
}
\lstset{style=mystyle}
\newcommand{\Sref}[1]{Section~\ref{#1}}
\newtheorem{hyp}{Hypothesis}


\title{Problem Set 4}
\date{Due: November 26, 2026}
\author{
	Applied Stats/Quant Methods 1\\
	Kanon Saito
}


\begin{document}
	\maketitle
	\section*{Instructions}
	\begin{itemize}
		\item Please show your work! You may lose points by simply writing in the answer. If the problem requires you to execute commands in \texttt{R}, please include the code you used to get your answers. Please also include the \texttt{.R} file that contains your code. If you are not sure if work needs to be shown for a particular problem, please ask.
		\item Your homework should be submitted electronically on GitHub.
		\item This problem set is due before 23:59 on Thursday November 26, 2026. No late assignments will be accepted.
	\end{itemize}



	\vspace{.5cm}
\section*{Question 1: Economics}
\vspace{.25cm}
\noindent 	
In this question, use the \texttt{prestige} dataset in the \texttt{car} library. First, run the following commands:

\begin{verbatim}
install.packages(car)
library(car)
data(Prestige)
help(Prestige)
\end{verbatim} 

\noindent We would like to study whether individuals with higher levels of income have more prestigious jobs. Moreover, we would like to study whether professionals have more prestigious jobs than blue and white collar workers.

\newpage
\begin{enumerate}
	
	\item [(a)]
	Create a new variable \texttt{professional} by recoding the variable \texttt{type} so that professionals are coded as $1$, and blue and white collar workers are coded as $0$ (Hint: \texttt{ifelse}).
	
	R:
	\lstinputlisting[linerange=39-39]{PS04_answers_KS.R}
	
	
	\item [(b)]
	Run a linear model with \texttt{prestige} as an outcome and \texttt{income}, \texttt{professional}, and the interaction of the two as predictors (Note: this is a continuous $\times$ dummy interaction.)
	
	R:
	\lstinputlisting[linerange=41-42]{PS04_answers_KS.R}
	
	\begin{verbatim}
	Call:
	lm(formula = prestige ~ income * professional, data = Prestige)
	
	Residuals:
	Min      1Q  Median      3Q     Max 
	-14.852  -5.332  -1.272   4.658  29.932 
	
	Coefficients:
	Estimate Std. Error t value Pr(>|t|)    
	(Intercept)         21.1422589  2.8044261   7.539 2.93e-11 ***
	income               0.0031709  0.0004993   6.351 7.55e-09 ***
	professional        37.7812800  4.2482744   8.893 4.14e-14 ***
	income:professional -0.0023257  0.0005675  -4.098 8.83e-05 ***
	---
	Signif. codes:  0 ‘***’ 0.001 ‘**’ 0.01 ‘*’ 0.05 ‘.’ 0.1 ‘ ’ 1
	
	Residual standard error: 8.012 on 94 degrees of freedom
	(4 observations deleted due to missingness)
	Multiple R-squared:  0.7872,	Adjusted R-squared:  0.7804 
	F-statistic: 115.9 on 3 and 94 DF,  p-value: < 2.2e-16
	\end{verbatim}
	\vspace{1cm}
	
	\item [(c)]
	Write the prediction equation based on the result.
	
	$\widehat{Prestige}=21.142+0.0032\times\text{income}+37.78\times\text{professional}-0.0023\times\text{income}\times\text{professional}$
	\vspace{1cm}

	\item [(d)]
	Interpret the coefficient for \texttt{income}.
	
	On average, for non-professionals (professional = 0), a one-unit increase in income is associated with a 0.0032-unit increase in job prestige.
	\vspace{1cm}	
	
	\item [(e)]
	Interpret the coefficient for \texttt{professional}.
	
	On average, when income is 0, being in a professional occupation (professional = 1) is associated with a 37.78-unit higher level of job prestige compared to non-professionals (professional = 0).
	\vspace{1cm}
	
	\item [(f)]
	What is the effect of a \$1,000 increase in income on prestige score for professional occupations? In other words, we are interested in the marginal effect of income when the variable \texttt{professional} takes the value of $1$. Calculate the change in $\hat{y}$ associated with a \$1,000 increase in income based on your answer for (c).
	
	R:
	\lstinputlisting[linerange=44-44]{PS04_answers_KS.R}
	
	The marginal effect of a one-unit increase of income on prestige:\\
	
	$\text{marginal effect of income}=\beta_1+\beta_3\times\text{professional}=0.0032-0.0023\times1=0.0009$\\
	
	Therefore, the effect of a \$1,000 increase in income on prestige is 0.9.
	\vspace{1cm}
	
	
	\item [(g)]
	What is the effect of changing one's occupation from non-professional to professional when her income is \$6,000? We are interested in the marginal effect of professional jobs when the variable \texttt{income} takes the value of $6,000$. Calculate the change in $\hat{y}$ based on your answer for (c).
	
	R: 
	\lstinputlisting[linerange=46-46]{PS04_answers_KS.R}
	
	The marginal effect of changing one's occupation from non-professional to professional on prestige score:\\
	
	$\text{marginal effect of professional}=\beta_2+\beta_3\times\text{income}=37.78-0.0023\times\text{income}$\\
	
	Therefore, when the income takes the value of $6,000$:\\
	$37.78-0.0023\times6000=23.98$
	
\end{enumerate}

\newpage

\section*{Question 2: Political Science}
\vspace{.25cm}
\noindent 	Researchers are interested in learning the effect of all of those yard signs on voting preferences.\footnote{Donald P. Green, Jonathan	S. Krasno, Alexander Coppock, Benjamin D. Farrer,	Brandon Lenoir, Joshua N. Zingher. 2016. ``The effects of lawn signs on vote outcomes: Results from four randomised field experiments.'' Electoral Studies 41: 143-150. } Working with a campaign in Fairfax County, Virginia, 131 precincts were randomly divided into a treatment and control group. In 30 precincts, signs were posted around the precinct that read, ``For Sale: Terry McAuliffe. Don't Sellout Virginia on November 5.'' \\

Below is the result of a regression with two variables and a constant.  The dependent variable is the proportion of the vote that went to McAuliffe's opponent Ken Cuccinelli. The first variable indicates whether a precinct was randomly assigned to have the sign against McAuliffe posted. The second variable indicates
a precinct that was adjacent to a precinct in the treatment group (since people in those precincts might be exposed to the signs).  \\

\vspace{.5cm}
\begin{table}[!htbp]
	\centering 
	\textbf{Impact of lawn signs on vote share}\\
	\begin{tabular}{@{\extracolsep{5pt}}lccc} 
		\\[-1.8ex] 
		\hline \\[-1.8ex]
		Precinct assigned lawn signs  (n=30)  & 0.042\\
		& (0.016) \\
		Precinct adjacent to lawn signs (n=76) & 0.042 \\
		&  (0.013) \\
		Constant  & 0.302\\
		& (0.011)
		\\
		\hline \\
	\end{tabular}\\
	\footnotesize{\textit{Notes:} $R^2$=0.094, N=131}
\end{table}

\vspace{.5cm}
\begin{enumerate}
	\item [(a)] Use the results from a linear regression to determine whether having these yard signs in a precinct affects vote share (e.g., conduct a hypothesis test with $\alpha = .05$).\\
	
	R:
	\lstinputlisting[linerange=49-51]{PS04_answers_KS.R}
	
	The p-value for the effect of having lawn signs on vote share is less than 0.05 ($p$ = 0.0097), which means it has a statistically significant effect on vote share.  Therefore, we can reject the null hypothesis that yard signs have no effect on vote share.
	
	\vspace{1cm}
	\item [(b)]  Use the results to determine whether being
	next to precincts with these yard signs affects vote
	share (e.g., conduct a hypothesis test with $\alpha = .05$).\\
	
	R:
	\lstinputlisting[linerange=54-56]{PS04_answers_KS.R}
	
	The p-value for the effect of being adjacent to a precinct assigned lawn signs on vote share is less than 0.05 ($p$ = 0.0016), which means it has a statistically significant effect on vote share.  Therefore, we can reject the null hypothesis that being adjacent to a precinct assigned lawn signs has no effect on vote share.
	\vspace{1cm}
	\item [(c)] Interpret the coefficient for the constant term substantively.\\
	
	The coefficient for the constant term indicates that, on average, when lawn signs are not assigned to the precinct ($X_1$) and not being adjacent to a precinct assigned lawn signs ($X_2$), the estimated vote share of Ken Cuccinelli is 0.302 (30.2\%).
	
	\vspace{1cm}
	
	\item [(d)] Evaluate the model fit for this regression.  What does this	tell us about the importance of yard signs versus other factors that are not modelled?\\
	
	The model fit is relatively low, with an R-squared value of 0.094. This indicates that having lawn signs in a precinct and being adjacent to a precinct assigned lawn signs together explain only 9.4\% of the variation in the vote share of Ken Cuccinelli. Therefore, other unmodeled factors (accounting for the remaining 90.6\% of the variation) play a much larger role in determining electoral outcomes than lawn signs do.
	
	
	
	
\end{enumerate}  


\end{document}
